主码筛选既需要暴露足够的 GPU 并行工作，又不能丢掉本应进入残差精排的候选。以查询为单位的扫描为每条查询启动一个 block；当 batch 很小时，即使每条查询访问大量向量，也只有少量 block 可供执行。GPU 上的选择操作本身就是一个重要的设计问题\cite{johnson_gpu}。v62 引入的分组扫描将已访问候选划分给多个 block，再合并各组保留的候选。它保证的是基于已计算主码分数的精确选择；IVF 路由和量化表示仍然具有近似性。

\subsection{候选划分与确定性顺序}

对固定查询，按探针顺序拼接选中的倒排列表，得到 $n_c$ 个候选出现位置。假设每个被访问数据库行只出现一次，搜索期间索引保持不变。对 $G=\gscan\in\{1,2,4,8\}$，block $(q,g)$ 负责区间
\begin{equation}
 U_g=\left[\left\lfloor\frac{g n_c}{G}\right\rfloor,
              \left\lfloor\frac{(g+1)n_c}{G}\right\rfloor\right),
 \qquad g=0,\ldots,G-1.
 \label{eq:5-partition}
\end{equation}
这里的区间表示相应候选，而不是它们的数据库 ID。设备端依据路由得到的候选数计算区间端点。各组候选数最多相差一个；一个较长的列表可以跨越组边界。这种划分平衡的是候选数量，而平均分配列表数量不一定能达到同样效果。每个线程通过探针前缀偏移找到自己的第一个列表，在扫描区间时向前移动列表游标。此阶段不读取残差码。

选择使用的 key 为
\begin{equation}
 K_q(x)=\bigl(\primarydist(q,x),p(x)\bigr),
 \qquad
 (a,p)\prec(b,r)\iff a<b\ \text{or}\ (a=b\ \text{and}\ p<r),
 \label{eq:5-primary-key}
\end{equation}
其中 $p(x)$ 是固定倒排存储中的物理行位置。它在该查询的候选中唯一，不依赖原子追加的先后顺序，也不是最终结果处理并列时使用的数据库 ID。当分数不含 NaN 时，这一定义为不同候选建立全序。局部比较和合并决策都使用相同顺序。要使选择不受划分方式影响，还要求各条比较路径为同一候选计算出相同的分数；排序规则本身不能消除浮点评分差异。

\subsection{精确局部保留与全局覆盖}

用 $\operatorname{Top}_c(V)$ 表示将 $V$ 按 $K_q$ 排序后最前面的 $\min(c,|V|)$ 个候选，其中 $c=\max\{k,\lceil\alpha k\rceil\}$。每个独立组保留完整的局部 top-$c$，而不是分得一个 $c/G$ 配额。这一区别不可省略：当 $c=4$、组数为二时，四个全局优胜候选可能全在同一组。如果每组只保留两个，即使总容量为四，仍会丢掉全局优胜候选。

\begin{proposition}[局部覆盖与截断合并]
\label{prop:5-local-coverage}
给定互不相交的组 $U_0,\ldots,U_{G-1}$、固定候选分数及共同全序，有
\begin{equation}
 \operatorname{Top}_c\!\left(\bigcup_g U_g\right)
 =\operatorname{Top}_c\!\left(\bigcup_g
                    \operatorname{Top}_c(U_g)\right).
 \label{eq:5-topc-union}
\end{equation}
对有序列表两两合并，并在每次合并后截取前 $c$ 项，保持这一结果；候选总数不足 $c$ 时也成立。
\end{proposition}
\begin{proof}
设 $x$ 同时属于全局 top-$c$ 和组 $U_g$。全局排在 $x$ 之前的候选少于 $c$ 个，因此组 $U_g$ 内排在它之前的候选也少于 $c$ 个。于是 $x\in\operatorname{Top}_c(U_g)$，所有局部结果的并集包含每个全局优胜候选。再从这个并集中取前 $c$，就精确恢复同一顺序下的全局结果。如果整个候选集合不足 $c$，所有成员均被保留。

对于合并，两个有序列表的前 $c$ 项就是二者并集的 top-$c$。被截去的候选在该并集中已经有至少 $c$ 个前驱，不可能再进入任何更大并集的 top-$c$。对合并树的各轮归纳即可得到结论。
\end{proof}

\begin{figure}[tbp]
\centering
% Schematic only: illustrative keys, not experimental observations.
\begin{tikzpicture}[
  font=\small\linespread{1}\selectfont,
  box/.style={draw=ThesisNavy!65,rounded corners=2pt,align=center,
    inner sep=5pt,text width=2.65cm,minimum height=2.05cm},
  merged/.style={box,text width=5.55cm,minimum height=1.05cm,
    fill=ThesisNavy!5},
  arr/.style={-{Latex[length=2mm]},draw=ThesisNavy!80,thick}]
\node[box] (g0) at (0,0) {$\operatorname{Top}_3(U_0)$\\[2pt]
  $\mathbf{(1,8)}$\\$\mathbf{(2,6)}$\\$(2,9)$};
\node[box] (g1) at (3.5,0) {$\operatorname{Top}_3(U_1)$\\[2pt]
  $(4,1)$\\$(5,4)$\\\textit{length 2}};
\node[box] (g2) at (7,0) {$\operatorname{Top}_3(U_2)$\\[2pt]
  $\varnothing$\\\textit{length 0}};
\node[box] (g3) at (10.5,0) {$\operatorname{Top}_3(U_3)$\\[2pt]
  $\mathbf{(2,3)}$\\$(7,2)$\\$(8,5)$};
\node[merged] (m0) at (1.75,-2.5) {Merge and retain first 3\\[2pt]
  $\mathbf{(1,8)},\ \mathbf{(2,6)},\ (2,9)$};
\node[merged] (m1) at (8.75,-2.5) {Merge and retain first 3\\[2pt]
  $\mathbf{(2,3)},\ (7,2),\ (8,5)$};
\node[merged,text width=6.8cm,fill=ThesisNavy!12] (out) at (5.25,-4.3)
  {Global primary top-3\\[2pt]
   $\mathbf{(1,8)},\ \mathbf{(2,3)},\ \mathbf{(2,6)}$};
\draw[arr] (g0.south) -- (m0.north west);
\draw[arr] (g1.south) -- (m0.north east);
\draw[arr] (g2.south) -- (m1.north west);
\draw[arr] (g3.south) -- (m1.north east);
\draw[arr] (m0.south) -- (out.north west);
\draw[arr] (m1.south) -- (out.north east);
\end{tikzpicture}

\caption{局部覆盖与截断式两两合并，示例中 $c=3$、$G=4$。条目是 $(\primarydist,p)$ key，属于示意而非测量数据。粗体条目是最终全局优胜候选。短组仅贡献有效前缀，空组不贡献条目。每次合并使用同一 key 顺序，并保留至多 $c$ 项。}
\label{fig:5-global-selection}
\end{figure}

CUDA 实现在共享内存中维护容量为 $C$ 的 key 对缓冲区。扫描 block 含 $T$ 个线程，$C$ 为二的幂且满足 $C\geq c+T$。压缩操作通过 bitonic 网络排序，保留至多 $c$ 个有效条目；当有效条目达到 $c$ 个时，将最后一个保留 key 记为局部阈值。在此之前，阈值是排在所有有效 key 之后的哨兵。两次压缩之间，仅追加排在阈值之前的候选。即使后续候选继续到来，阈值对应的 $c$ 个见证候选仍使它成为已扫描前缀中最优 $c$ 项的安全上界。

每个 chunk 最多追加 $T$ 项。处理完一个 chunk 后，如果下一个 chunk 可能使缓冲区越界，就先压缩。压缩后至多剩下 $c$ 项，$C\geq c+T$ 保证容得下下一个 chunk；不压缩时，显式容量检查提供同样保证。该不变量防止溢出，也不限制一个线程可以遇到多少个有竞争力的候选。固定容量的逐线程保留缓冲区并不具有相同的覆盖保证。

另一种方案是先将全部 $n_c$ 个分数／位置对写出，再调用独立选择器。这保持相同目标，但增加一个 $O(n_c)$ 中间数组及其写入。这里将打包主码的评分直接送入有界局部缓冲区，以重复局部排序换取 $O(Gc)$ 规模的中间输出。基于寄存器的选择也是已有的 GPU 设计\cite{johnson_gpu}；本轮实验评估选定的扫描变体，并未证明其优于所有库选择器。

\subsection{合并实现、存储与边界情况}

当 $G>1$ 时，每条查询由一个 256 线程的 block 执行 $\log_2 G$ 轮两两合并。有效长度为 $a,b$ 的一对列表输出 $\min(c,a+b)$ 个 key。Merge Path 将输出划分为连续的线程片段：先二分查找每个片段对应的输入偏移，再由该线程顺序合并自己的片段\cite{gpu_merge_path}。所有比较都遵循\cref{eq:5-primary-key}，各轮之间通过 block 屏障同步。\Cref{fig:5-global-selection} 展示了截断不变量。除划分查找与同步工作外，合并树全部内部节点合计最多输出 $c(G-1)$ 个 key 对。

两个全局分数／位置工作区的容量分别为 $[b_{\rm cap},G,c]$ 和 $[b_{\rm cap},G/2,c]$，在各轮间交替使用。分数和位置各占四字节，因此二者的总载荷为
\begin{equation}
 W_{\rm pairs}=8b_{\rm cap}c\left(G+\frac{G}{2}\right)
 \quad\text{bytes},\qquad G>1.
 \label{eq:5-group-workspace}
\end{equation}
各组计数还需要 $4b_{\rm cap}G$ 字节；配置共享界存储时另需 $8b_{\rm cap}$ 字节。这些空间位于最终保留数组及其他搜索工作区之外。分配容量在 graph capture 之前增长，之后保留复用，因此以后选择较小的 $G$ 不一定释放预留显存。合并只在共享内存中保存列表长度，而不存放全部 $Gc$ 个 key。

局部有效前缀由计数界定，而不是仅依靠哨兵分数。压缩时，排序缓冲区填充 $(+\infty,\mathrm{INT\_MAX})$；有效位置必须小于这个保留位置。空组计数为零。最终保留数组将未使用槽位填为 $(+\infty,-1)$，残差精排跳过负位置。这两种哨兵用途不同：负位置填充项不能作为有效 key 参与合并。当 $G=1$ 时，直接写最终保留数组，无需合并。

较早的查询级路径由实现标签 $G=0$ 选择，表示每条查询一个 block，而不是零个 block。其默认主码选择按分数精确，但仅比较距离的截断规则可能在并列处保留与完整 key 路径不同的成员。有损的 \texttt{K\_LOCAL=4} 替代路径需要在允许该路径的构建中通过 \texttt{JHQ\_EXACT\_TOPCK=0} 开启；paper-faithful 构建强制使用精确选择。它不是被测默认路径。\Cref{sec:7-4} 报告了 64 组分组扫描比较：主码选择没有非并列差异，在实际使用查询上最终输出一致。这些检查支持覆盖范围内的实现行为；本命题证明的是有条件的选择性质，不是原空间精确 kNN，也不是与所有旧并列规则无条件等价。
